\emph{Setting.} The algebraic limit $L_{p,q}$ is the union of the floors under the inclusions of generators (\claimref{obj.tower.tower}), that is, the Clifford algebra of the countable-dimensional quadratic space $V \oplus H^{\infty}$ spanned by the ground and plane generators; an assignment of elements to the generators that satisfies the Clifford relations (each image squares to the square of its generator, distinct images anticommute) therefore extends uniquely to an algebra homomorphism (\claimref{std-algebra.existence-uniqueness}), floor by floor and hence on the union.

\emph{(i).} Write $Q(g)$ for the square of a generator $g$. Every image is a product of two anticommuting unit generators, except $\Phi(e_{1}) = e_{1}$. For two anticommuting generators $g, h$ of the target, $(gh)^{2} = -g^{2}h^{2}$, so $\Phi(x_{i})^{2} = -y_{i}^{2} = x_{i}^{2}$, $\Phi(e_{j})^{2} = (e_{1}f_{j-1})^{2} = -(1)(-1) = 1$ for $j \ge 2$, and $\Phi(f_{j})^{2} = (e_{1}e_{j+1})^{2} = -1$. For anticommutation: two images $e_{1}g$, $e_{1}h$ with $g \ne h$ distinct from $e_{1}$ satisfy $(e_{1}g)(e_{1}h) = -gh$, whose symmetrisation vanishes; $e_{1}$ anticommutes with every $e_{1}g$, since $e_{1}(e_{1}g) = g = -(e_{1}g)e_{1}$; and $\Phi(x_{i}) = y_{i}e_{1}$ against $e_{1}h$ or $e_{1}$ is the same computation with $y_{i}$ in place of $g$, the right-hand placement of $e_{1}$ changing only an overall sign. The images are pairwise distinct products (the map $g \mapsto$ the second factor is injective onto $\{y_{i}\} \cup \{f_{j-1}\} \cup \{e_{j+1}\}$), so $\Phi$ is an algebra homomorphism $L_{p,q} \to L_{q,p}$. The same formula read from $(q,p)$ to $(p,q)$ defines $\Psi \colon L_{q,p} \to L_{p,q}$, and on generators $\Psi\Phi(x_{i}) = \Psi(y_{i})\Psi(e_{1}) = x_{i}e_{1}e_{1} = x_{i}$, $\Psi\Phi(e_{j}) = e_{1}\cdot e_{1}e_{j} = e_{j}$ ($j \ge 2$), $\Psi\Phi(f_{j}) = e_{1}\cdot e_{1}f_{j} = f_{j}$, and symmetrically $\Phi\Psi = \mathrm{id}$; so $\Phi$ is an isomorphism. The images of plane generators involve plane generators only, so $\Phi(B_{\infty}) \subseteq B_{\infty}$ and likewise $\Psi(B_{\infty}) \subseteq B_{\infty}$, whence $\Phi(B_{\infty}) = B_{\infty}$. The images of the generators of $T_{k}$ use plane indices up to $k+1$, which gives the floor statements; $\Phi(T_{k}) \not\subseteq T'_{k}$ because $\Phi(f_{k}) = e_{1}e_{k+1}$, and $\dim \Phi(T_{k}) = \dim T'_{k} < \dim T'_{k+1}$, so no floor goes onto a floor. Finally $\Phi(x_{i}x_{j}) = y_{i}e_{1}y_{j}e_{1} = -y_{i}y_{j}$. \emph{Completion.} For the frame $*$-structure of Theorem \claimref{tower.structure-tower-c-norms}, $g^{*} = g^{-1}$ for every generator. Each $\Phi(g)$ is a product of frame generators, hence unitary, so $\Phi(g)^{*} = \Phi(g)^{-1} = \Phi(g^{-1}) = \Phi(g^{*})$; $\Phi \circ *$ and $* \circ \Phi$ are conjugate-linear antihomomorphisms (after complexification) agreeing on generators, so they agree and $\Phi$ is a $*$-isomorphism. Its restriction $T_{k,\C} \to T'_{k+1,\C}$ is an injective $*$-homomorphism of finite-dimensional C*-algebras, hence isometric for their unique C*-norms; so $\Phi$ and $\Psi$ are isometric on the union and extend to mutually inverse isomorphisms of the completions (for $p - q$ even, the completion of Theorem \claimref{tower.car-limit-even-charge}), carrying the closure of $B_{\infty}$ onto itself.

\emph{(ii).} If $n \ge 1$ then $x_{1} \notin B_{\infty}$, while for $n' = 0$ the plane algebra is the whole limit, so no pair isomorphism exists between $n \ge 1$ and $n' = 0$. For $n, n' \ge 1$: if the pairs are isomorphic, so are the relative commutants of $B_{\infty}$, which are $\Cl^{0}(p,q)$ and $\Cl^{0}(p',q')$ by Theorem \claimref{synthesis.plane-commutant-even-ground}(i); by its part (ii) these determine $n$ and the class $\{r, -r\} \bmod 8$. Conversely, (i) gives $(p,q) \to (q,p)$. For the shift $(p,q) \to (p-4, q+4)$ when $p \ge 4$: take four ground generators $x_{1}, \dots, x_{4}$ of square $+1$ and $\omega_{4} = x_{1}x_{2}x_{3}x_{4}$, so $\omega_{4}^{2} = 1$, $\omega_{4}$ anticommutes with $x_{1}, \dots, x_{4}$ and commutes with every other generator. The elements $z_{i} = x_{i}\omega_{4}$ ($i \le 4$) satisfy $z_{i}^{2} = -1$, anticommute pairwise and with every other generator, and $z_{1}z_{2}z_{3}z_{4} = \omega_{4}$, so $x_{i} = z_{i}\omega_{4}$; together with the remaining ground generators and all plane generators, unchanged, they present $L_{p,q}$ as $L_{p-4,q+4}$ with $B_{\infty}$ fixed pointwise. The case $q \ge 4$ is the same with squares $-1$. For fixed $n$, $r = 2p - n$, so $r' \equiv r \pmod 8$ means $p' \equiv p \pmod 4$, reached by shifts within $0 \le p \le n$, and $r' \equiv -r$ means $p' \equiv n - p \pmod 4$, reached by one reversal followed by shifts.

\emph{(iii).} Over $\R$: for a blade $B$ and a generator $g$ (the blade argument of Theorem \claimref{synthesis.plane-commutant-even-ground}, Step 1): if $g \notin B$ then $gB = (-1)^{|B|}Bg$, which is supercommutation; if $g \in B$ then $gB = (-1)^{|B|-1}Bg$, which is not. Conjugation by $g$ acts diagonally on blades, so the supercommutant of all plane generators is spanned by the blades containing no plane generator, which span $\Cl(p,q)$. An isomorphism of graded pairs restricts to an isomorphism of these supercommutants, and $\Cl(p,q) \not\cong \Cl(q,p)$ for $r \not\equiv 0, 4 \pmod 8$ (the types for $r$ and $-r$ differ: $\R \oplus \R$ against $\C$, $\R$ against $\HH$, $\C$ against $\HH \oplus \HH$; Theorem \claimref{transport.classification-real-clifford-algebras}). For $\Phi\alpha\Phi^{-1}$, evaluate on the generators of $L_{q,p}$ through $\Psi$: $y_{i} = \Phi(x_{i})\Phi(e_{1}) \mapsto (-\Phi(x_{i}))(-\Phi(e_{1})) = y_{i}$; $e_{1} \mapsto -e_{1}$; $f_{j} = \Phi(e_{1})\Phi(e_{j+1}) \mapsto f_{j}$; $e_{j} = \Phi(e_{1})\Phi(f_{j-1}) \mapsto e_{j}$ for $j \ge 2$. So $\Phi\alpha\Phi^{-1}$ negates $e_{1}$ alone, which is $\operatorname{Ad}(e_{1}) \circ \alpha'$, and $e_{1}\alpha'(e_{1}) = -e_{1}^{2} = -1$. Its supercommutant of the planes is $\Phi(\Cl(p,q))$, generated by the $y_{i}e_{1}$. The blade argument holds verbatim after complexification, giving $\Cl(p,q)_{\C} \cong \Cl(p+q, 0)_{\C}$, which depends on $n$ alone; so over $\C$ the graded pair does not separate $r$ from $-r$ either.

\emph{(iv).} With $p - q \equiv 1 \pmod 4$ the repairs have residues $r \pm 1$, both even; $r + 1 \equiv -(r - 1) \pmod 8$ would need $2r \equiv 0 \pmod 8$, impossible for $r$ odd. By (ii) the pairs are not isomorphic. Concretely the planes' commutant is $\Cl^{0}$ of the repair, whose centre is spanned by $1$ and the volume element $\omega$ with $\omega^{2} = (-1)^{r'/2}$ (Theorem \claimref{synthesis.plane-commutant-even-ground}, Step 3): $r' = r + 1 \equiv 2 \pmod 4$ gives $\omega^{2} = -1$, a centre $\C$, and $r' = r - 1 \equiv 0 \pmod 4$ gives $\omega^{2} = 1$, a centre $\R \oplus \R$. The fused algebras themselves are isomorphic by Fact \claimref{rule2.imported-rule-2-matrix-units-forgotten}.

The check verifies (i) on truncations ($n \le 5$, two floors), including the inverse and the transported grading, the shift of (ii), and the centres of (iv), in exact arithmetic.
